\section{Preliminaries}\label{sec:preliminaries}
In this section, we introduce some preliminaries that will be used in our later analysis and discussion.

\subsection{Reviews on diffusion-based Monte Carlo}
We consider the following stochastic differential equation (SDE) on $\mathbb{R}^d$:
\begin{equation}\label{eqn:sde}
\d X_t=-\nabla V(X_t)\d t+ \sigma \d B_t,
\end{equation}
where $V: \mathbb{R}^d\rightarrow \mathbb{R}$ is a potential function, $\sigma$ is the constant noise intensity, and $B_t$ is a $w$-dimensional standard Brownian motion. We commence by establishing the necessary assumptions regarding the potential $V$ to ensure the well-posedness of the underlying dynamics.
\begin{assumption}\label{assum:SDE}
    \begin{enumerate}
        \item There exists a positive constant $K_1$ such that
        \[
            -2x\cdot \nabla V(x) \leq K_1(|x|^2 + 1) \quad \text{for all } x\in\mathbb{R}^d,
        \]
        and $V$ is locally uniformly continuous.\label{subassu:K1}
        \item There exists a positive constant $\delta_0$ such that for all $R>0$ and $x,z\in\mathbb{R}^d$ satisfying $|x|\leq R$, $|z|\leq R$, and $|x-z|\leq \delta_0$, the inequality
        \[
            2(x-z)\cdot(\nabla V(z)-\nabla V(x))  \leq K_R|x-z|^2
        \]
        holds, where $K_R$ is a positive constant dependent on $R$.
        \item There exist a positive constant $\alpha$, a compact set $C\subset\mathbb{R}^d$, a measurable function $f:\mathbb{R}^d\mapsto[1,\infty)$, and a twice continuously differentiable function $U:\mathbb{R}^d\mapsto\mathbb{R}_+$ satisfying
        \[
            \mathcal{L}U(x)\leq -\alpha f(x) + 1_C(x) \quad \text{for all } x\in\mathbb{R}^d.
        \]
    \end{enumerate}
\end{assumption}

Under Assumption \ref{assum:SDE}, the existence and uniqueness of the solution to the It\^{o} SDE \eqref{eqn:sde} are theoretically guaranteed. Furthermore, the process is ergodic and admits a unique invariant distribution $p_\infty$; see, e.g., \cite[Theorem 2.2, Theorem 2.4, Remark 5.5, Lemma 6.4]{xi2019jump}. The probability density function of the solution process $X_t$ at time $t$ and at position $x$, denoted by $p_t(x)$, evolves according to the Fokker--Planck equation
\begin{equation}
\label{eqn::fp}
   \partial_t p_t(x)  = \nabla\cdot\big(\nabla V(x) p_t(x)\big) + \frac{\sigma^2}{2} \triangle p_t(x).
\end{equation}
The invariant distribution $p_\infty$ is defined as the long-time limit of the transient density $p_t$, namely, $p_\infty:=\underset{t\to +\infty}{\lim}p_t$. The invariant measure of the SDE \eqref{eqn:sde} admits the explicit analytical form
\begin{equation}\label{p_infty}
    p_\infty(x)= Z^{-1}\exp\left\{-\frac{2V(x)}{\sigma^2}\right\},
\end{equation}
which is widely recognized in statistical mechanics as the Boltzmann distribution. Here, $Z$ is the normalization constant---commonly referred to as the partition function---ensuring that $p_\infty$ integrates to unity.


The ergodicity of the solution process $X_t$ implies that the invariant distribution $p_\infty$ can be accurately approximated by the time average along a sufficiently long trajectory:
\begin{equation}\label{eqn:ergodicity}
    p_\infty(x) \approx \frac{1}{T}\int_0^T\delta(x-X_t)\d t.
\end{equation}
This ergodic property underpins Monte Carlo methods, such as the Euler--Maruyama scheme, facilitating the numerical sampling of the invariant distribution. However, in the weakly stochastic regime (i.e., when the noise intensity $\sigma$ is small) and when the potential $V$ features multiple deep wells separated by high energy barriers, the simulation time $T$ required in \eqref{eqn:ergodicity} becomes prohibitively large. In such scenarios, the process $X_t$ exhibits metastability, remaining trapped within local minima for exponentially long durations before transitioning to other wells---a phenomenon characteristic of rare events. This computational inefficiency arises from the continuous, diffusive nature of Brownian motion, which yields an average spatial displacement of order $\mathcal{O}(\sqrt{\Delta t})$ over a small time step $\Delta t$, thereby fundamentally limiting the exploration of the state space.





\subsection{Generators and Dirichlet forms for diffusions and jump-diffusions}\label{subsubsec:gen_dirichlet}
%========================================================

We collect the infinitesimal generators and the associated energy (Dirichlet) forms for
(i) the overdamped Langevin diffusion and
(ii) the jump--diffusions.
Throughout, we work on the Hilbert space $L^2(p_\infty)$ with inner product
\[
\langle f,g\rangle_{p_\infty}:=\int_{\mathbb{R}^d} f(x)\,g(x)\,p_\infty(x)\,\mathrm{d}x,
\]
and use $C_c^\infty(\mathbb{R}^d)$ as a core of test functions.

%--------------------------------------------------------
\paragraph{1.Diffusion.}
The (backward) generator $\mathcal{L}$ of diffusion SDE \eqref{eqn:sde} acts on $\varphi\in C_c^\infty(\mathbb{R}^d)$ as
\begin{equation}\label{eq:generator_diffusion}
(\mathcal{L}\varphi)(x)= -\nabla V(x)\cdot\nabla\varphi(x)+\frac{\sigma^2}{2}\Delta\varphi(x).
\end{equation}
\begin{definition}[Dirichlet form]
    The associated Dirichlet form of diffusion SDE \eqref{eqn:sde} is defined by
\begin{equation}\label{eq:dirichlet_diffusion}
\mathcal{E}_{\mathrm{diff}}(f,g):=-\langle f,\mathcal{L}g\rangle_{p_\infty}
=\frac{\sigma^2}{2}\int_{\mathbb{R}^d}\nabla f(x)\cdot \nabla g(x)\,p_\infty(x)\,\mathrm{d}x,
\qquad f,g\in C_c^\infty(\mathbb{R}^d),
\end{equation}
so that
\[
\mathcal{E}_{\mathrm{diff}}(f,f)=\frac{\sigma^2}{2}\int_{\mathbb{R}^d}|\nabla f(x)|^2\,p_\infty(x)\,\mathrm{d}x\ge 0.
\]
\end{definition}
Equivalently,
\begin{equation}\label{eq:generator_divergence_form}
\mathcal{L}g=\frac{\sigma^2}{2}\,p_\infty^{-1}\nabla\cdot\big(p_\infty\nabla g\big),
\end{equation}
which makes the $L^2(p_\infty)$-symmetry of $\mathcal{E}_{\mathrm{diff}}$ immediately.

%--------------------------------------------------------
\paragraph{2. Jump--diffusion.}
Consider the autonomous L\'evy SDE 
\begin{equation}\label{eqn:sdegeneral}
\mathrm{d}Z_t=b(Z_{t-})\,\mathrm{d}t+\sigma\,\mathrm{d}B_t+\mathrm{d}L_t,
\end{equation}
where $L_t$ is a pure-jump L\'evy process with L\'evy measure $\nu(\mathrm{d}r)$ \cite{applebaum2009levy}.
We adopt the (forward) jump-operator
\[
\int_{\mathbb{R}^d\setminus\{0\}}\big(p(x-r)-p(x)\big)\,\nu(\mathrm{d}r),
\]
namely the backward jump generator
\[
(\mathcal{J}\varphi)(x):=\int_{\mathbb{R}^d\setminus\{0\}}\big(\varphi(x+r)-\varphi(x)\big)\,\nu(\mathrm{d}r),
\qquad \varphi\in C_c^\infty(\mathbb{R}^d),
\]
whenever the integral is well-defined. Hence the full generator $\mathcal{A}$ is
\begin{equation}\label{eq:generator_levy}
(\mathcal{A}\varphi)(x)
=b(x)\cdot\nabla\varphi(x)+\frac{\sigma^2}{2}\Delta\varphi(x)
+\int_{\mathbb{R}^d\setminus\{0\}}\big(\varphi(x+r)-\varphi(x)\big)\,\nu(\mathrm{d}r).
\end{equation}

\begin{definition}[Dirichlet form]
Suppose $p_\infty$ is invariant for \eqref{eqn:sdegeneral}, we define the Dirichlet form of jump-diffusion SDE \eqref{eqn:sdegeneral} as
\[
\mathcal{E}_{\mathrm{levy}}(f,g):=-\langle f,\mathcal{A}g\rangle_{p_\infty},
\qquad f,g\in C_c^\infty(\mathbb{R}^d),
\]
whenever the integrals are well-defined.
\end{definition}\label{def:dirichletlevy}

%--------------------------------------------------------

\subsection{Mixing, spectral gap, and mixing-time comparison}\label{subsubsec:L2_gap_mixing}
%========================================================

In this section, we outline the theoretical framework for analyzing the \emph{mixing} properties of the sampled processes---which characterizes the convergence of their probability distributions---through the lens of test functions and the spectral gap of the associated Markov semigroup. Specifically, we adopt an $L^2(p_\infty)$-based criterion for mixing, which is equivalent to convergence in the $\chi^2$-divergence for probability densities, and connect it to the corresponding Poincar\'e constant. This formulation establishes a rigorous comparison principle between the standard diffusion sampler \eqref{eqn:sde} and the proposed L\'evy-score sampler \eqref{eqn:sdegeneral}. Provided both processes admit the identical invariant measure $p_\infty$, comparing their convergence rates mathematically reduces to evaluating their respective $L^2(p_\infty)$ spectral gaps.

\begin{definition}[Semigroups, and invariance]
    Let $(P_t^{\mathrm{diff}})_{t\ge 0}$ and $(P_t^{\mathrm{levy}})_{t\ge 0}$ be the Markov semigroups induced by
\eqref{eqn:sde} and \eqref{eqn:sdegeneral}, respectively:
\[
(P_t^{\mathrm{diff}}f)(x):=\mathbb{E}[f(X_t)\mid X_0=x],
\qquad
(P_t^{\mathrm{levy}}f)(x):=\mathbb{E}[f(Z_t)\mid Z_0=x].
\]
\end{definition}

We assume that $p_\infty$ is invariant for both dynamics \eqref{eqn:sde} and \eqref{eqn:sdegeneral}; in particular,
\begin{equation}\label{eq:pinfty_invariant}
p_\infty(P_t f)=p_\infty(f),\qquad t\ge 0,
\end{equation}
for all bounded measurable $f$ for which the expectations are well-defined.

Throughout, we work on $L^2(p_\infty)$ with inner product
\[
\langle f,g\rangle_{p_\infty}:=\int_{\R^d} f(x)\,g(x)\,p_\infty(x)\,\d x,
\qquad
\|f\|_{L^2(p_\infty)}^2:=\langle f,f\rangle_{p_\infty},
\]
and define
\[
p_\infty(f):=\int_{\R^d} f(x)\,p_\infty(x)\,\d x,
\qquad
\mathrm{Var}_{p_\infty}(f):=\|f-p_\infty(f)\|_{L^2(p_\infty)}^2.
\]

%--------------------------------------------------------
\paragraph{From distributional distance to test functions (why ``test functions'' appear).}
A difference of probability measures $\mu-\nu$ is a linear functional acting on test functions:
\[
(\mu-\nu)(f)=\int f\,\d\mu-\int f\,\d\nu.
\]
Many common metrics can be written as a supremum over a class of test functions; for example,
$\|\mu-\nu\|_{\mathrm{TV}}=\sup_{\|f\|_\infty\le 1}|(\mu-\nu)(f)|$ and
$W_1(\mu,\nu)=\sup_{\mathrm{Lip}(f)\le 1}|(\mu-\nu)(f)|$.
In this paper we adopt an $L^2(p_\infty)$-based notion: if $\mu\ll p_\infty$ and we write
$h:=\frac{\d\mu}{\d p_\infty}$, then the $\chi^2$-divergence is
\[
\chi^2(\mu\|p_\infty)=\int (h-1)^2\,\d p_\infty=\|h-1\|_{L^2(p_\infty)}^2.
\]
Thus controlling $\|h_t-1\|_{L^2(p_\infty)}$ is a concrete way to quantify convergence of the law at time $t$ to
$p_\infty$, provided the evolving law stays absolutely continuous w.r.t.\ $p_\infty$. The spectral-gap method is
precisely an $L^2$ contraction argument for a semigroup, and therefore naturally expressed in terms of test functions.

%--------------------------------------------------------
\begin{definition}[$L^2$ energy identity]
Let $(P_t)_{t\ge 0}$ denote either $(P_t^{\mathrm{diff}})$ or $(P_t^{\mathrm{levy}})$, with generator
$\mathcal{G}\in\{\mathcal{L},\mathcal{A}\}$ acting on the core $C_c^\infty(\R^d)$ introduced in
Section~\ref{subsubsec:gen_dirichlet}. For $f\in C_c^\infty(\R^d)$ define the centered function
\[
g_t:=P_t f-p_\infty(f).
\]
\end{definition}
\begin{proposition}
    The centered function $g_t\in L^2(p_\infty)$ and it can be shown that
\begin{equation}\label{eq:L2_energy_identity}
\frac{\d}{\d t}\|g_t\|_{L^2(p_\infty)}^2
=
2\,\mathrm{Re}\,\langle g_t,\mathcal{G}g_t\rangle_{p_\infty}
=
-2\,\mathrm{Re}\,\langle g_t,-\mathcal{G}g_t\rangle_{p_\infty}.
\end{equation}
\end{proposition}

\begin{proof}
For any $f\in C_c^\infty(\R^d)$ and set $g_t:=P_t f-p_\infty(f)$. Since $p_\infty(f)$ is constant in $t$ and $P_t$
satisfies the Kolmogorov backward equation on the core, we have
$\partial_t g_t=\partial_t(P_t f)=\mathcal{G}P_t f$. Moreover, $\mathcal{G}$ annihilates constants, hence
$\mathcal{G}g_t=\mathcal{G}P_t f$, and therefore $\partial_t g_t=\mathcal{G}g_t$. Differentiating the squared norm gives
\[
\frac{\d}{\d t}\|g_t\|_{L^2(p_\infty)}^2
=
\frac{\d}{\d t}\langle g_t,g_t\rangle_{p_\infty}
=
\langle \partial_t g_t,g_t\rangle_{p_\infty}+\langle g_t,\partial_t g_t\rangle_{p_\infty}
=
\langle \mathcal{G}g_t,g_t\rangle_{p_\infty}+\langle g_t,\mathcal{G}g_t\rangle_{p_\infty}.
\]
Taking real parts yields
$\frac{\d}{\d t}\|g_t\|_{L^2(p_\infty)}^2
=2\,\mathrm{Re}\,\langle g_t,\mathcal{G}g_t\rangle_{p_\infty}$, as claimed.
\end{proof}

%--------------------------------------------------------
\begin{definition}[Dissipative quadratic form (coercivity)]
For a (possibly non-reversible) Markov semigroup with generator $\mathcal G$ on the core $C_c^\infty(\R^d)$,
we define the \emph{dissipative quadratic form}
\begin{equation}\label{eq:dissipative_form_general}
\mathcal{E}_{\mathcal G}(f,f)
:=
-\mathrm{Re}\,\langle f,\mathcal G f\rangle_{p_\infty},
\qquad f\in C_c^\infty(\R^d).
\end{equation}
\end{definition}
When $\mathcal G$ is $L^2(p_\infty)$-reversible, $\mathcal{E}_{\mathcal G}(f,f)=-\langle f,\mathcal G f\rangle_{p_\infty}$
is the usual (symmetric) Dirichlet form.

\begin{definition}[$L^2(p_\infty)$ coercivity (Poincar\'e) constant]
We define
\begin{equation}\label{eq:coercivity_def}
\lambda(\mathcal G)
:=
\inf\Big\{
\frac{\mathcal E_{\mathcal G}(h,h)}{\|h\|_{L^2(p_\infty)}^2}:\ 
h\in C_c^\infty(\R^d),\ p_\infty(h)=0,\ h\neq 0
\Big\}.
\end{equation}
\end{definition}\label{def:coercivity}
Equivalently, $\lambda(\mathcal G)>0$ is the largest constant such that
\begin{equation}\label{eq:coercivity_inequality}
\mathcal E_{\mathcal G}(h,h)\ \ge\ \lambda(\mathcal G)\,\|h\|_{L^2(p_\infty)}^2,
\qquad h\in C_c^\infty(\R^d),\ p_\infty(h)=0.
\end{equation}
In the reversible case, $\lambda(\mathcal G)$ coincides with the standard $L^2$ spectral gap.

\begin{definition}[Poincar\'e inequality]
By Definition \ref{def:coercivity}, $\lambda(\mathcal{G})>0$ is the largest constant such that
\begin{equation}\label{eq:poincare_inequality}
\mathcal{E}(h,h)\ \ge\ \lambda(\mathcal{G})\,\mathrm{Var}_{p_\infty}(h),
\qquad h\in C_c^\infty(\R^d),
\end{equation}
holds.
\end{definition}

%--------------------------------------------------------
\begin{proposition}[Exponential $L^2$ convergence from a spectral gap]
Assume that \eqref{eq:poincare_inequality} holds with $\lambda(\mathcal{G})>0$. Then for all $f\in C_c^\infty(\R^d)$,
\begin{equation}\label{eq:L2_mixing_bound}
\|P_t f-p_\infty(f)\|_{L^2(p_\infty)}
\le
e^{-\lambda(\mathcal{G})t}\,\|f-p_\infty(f)\|_{L^2(p_\infty)},
\qquad t\ge 0.
\end{equation}
\end{proposition}

\begin{proof}
Let $g_t=P_t f-p_\infty(f)$ as above. From \eqref{eq:L2_energy_identity} and \eqref{eq:dirichlet_form_general},
\[
\frac{\d}{\d t}\|g_t\|_{L^2(p_\infty)}^2
=
-2\,\mathrm{Re}\,\mathcal{E}(g_t,g_t).
\]
Since $\mathcal{E}(h,h)\ge 0$ for all $h$ and $\mathcal{E}(h,h)\in\R$ for the Dirichlet forms considered here, this is
\[
\frac{\d}{\d t}\|g_t\|_{L^2(p_\infty)}^2
=
-2\,\mathcal{E}(g_t,g_t).
\]
Applying the Poincar\'e inequality \eqref{eq:poincare_inequality} to $h=g_t$ yields
\[
\frac{\d}{\d t}\|g_t\|_{L^2(p_\infty)}^2
\le
-2\,\lambda(\mathcal{G})\,\mathrm{Var}_{p_\infty}(g_t).
\]
By invariance \eqref{eq:pinfty_invariant}, $p_\infty(g_t)=p_\infty(P_t f)-p_\infty(f)=0$, hence
$\mathrm{Var}_{p_\infty}(g_t)=\|g_t\|_{L^2(p_\infty)}^2$. Therefore
\[
\frac{\d}{\d t}\|g_t\|_{L^2(p_\infty)}^2\le -2\,\lambda(\mathcal{G})\,\|g_t\|_{L^2(p_\infty)}^2.
\]
Gronwall's inequality implies
$\|g_t\|_{L^2(p_\infty)}^2\le e^{-2\lambda(\mathcal{G})t}\|g_0\|_{L^2(p_\infty)}^2$,
which is equivalent to \eqref{eq:L2_mixing_bound}.
\end{proof}

%--------------------------------------------------------
\begin{definition}[$L^2$ mixing time]
For any $\delta\in(0,1)$, define the \emph{$L^2(p_\infty)$ mixing time} of $(P_t)$ by
\begin{equation}\label{eq:L2_mixing_time}
t_{\mathrm{mix}}^{(2)}(\delta)
:=
\inf\Big\{t\ge 0:\ 
\|P_t f - p_\infty(f)\|_{L^2(p_\infty)}
\le
\delta\,\|f-p_\infty(f)\|_{L^2(p_\infty)}\ \ \forall f\in L^2(p_\infty)\Big\}.
\end{equation}
\end{definition}
The estimate \eqref{eq:L2_mixing_bound} yields the standard bound
\begin{equation}\label{eq:tmix_gap_bound}
t_{\mathrm{mix}}^{(2)}(\delta)\ \le\ \frac{1}{\lambda(\mathcal{G})}\log\frac{1}{\delta}.
\end{equation}

Now we are ready for convenient comparison principles between
the diffusion sampler \eqref{eqn:sde} and the L\'evy-score sampler 
\eqref{eqn:sdegeneral} (when they share the same invariant measure $p_\infty$).
\begin{enumerate}
\item \textbf{(Connection to $\chi^2$-mixing of distributions.)}
If an initial law $\mu$ is absolutely continuous with respect to $p_\infty$, write
\(h = \frac{d\mu}{d p_\infty} \in L^2(p_\infty)\), and let \(\mu_t := \mu P_t\) denote
the law at time~\(t\). Under mild regularity assumptions ensuring that
\(\mu_t \ll p_\infty\) and that the adjoint semigroup \(P_t^\ast\) acts on
\(L^2(p_\infty)\), one has
\[
\frac{d\mu_t}{d p_\infty} = P_t^\ast h =: h_t.
\]
Since \(\chi^2(\mu_t\|p_\infty) = \|h_t - 1\|_{L^2(p_\infty)}^2\), this implies
\[
\chi^2(\mu_t\|p_\infty)
=
\|h_t - 1\|_{L^2(p_\infty)}^2
=
\|P_t^\ast(h-1)\|_{L^2(p_\infty)}^2.
\]
Note that \(h-1\) has mean zero with respect to \(p_\infty\). The $L^2(p_\infty)$
spectral gap bound \eqref{eq:L2_mixing_bound} provides exponential contraction
of \(P_t\) on the mean-zero subspace of \(L^2(p_\infty)\); moreover, \(P_t\) and
its adjoint \(P_t^\ast\) have the same $L^2(p_\infty)$ operator norm, so the
same exponential contraction estimate holds for \(P_t^\ast\). Applying this to
the mean-zero function \(h-1\) yields
\[
\chi^2(\mu_t\|p_\infty)^{1/2}
=
\|P_t^\ast(h-1)\|_{L^2(p_\infty)}
\le
e^{-\lambda(\mathcal{G})t}\,\|h-1\|_{L^2(p_\infty)}
=
e^{-\lambda(\mathcal{G})t}\,\chi^2(\mu\|p_\infty)^{1/2}.
\]
Thus the $L^2$-mixing time \(t_{\mathrm{mix}}^{(2)}\) corresponds to exponential
decay of \(\chi^2(\mu_t\|p_\infty)\) for densities.

\item \textbf{(Mixing-time comparison principle.)}
At the level of the $L^2(p_\infty)$ contraction bound, comparing the convergence speed of the two samplers reduces to
comparing their spectral gaps:
\[
\lambda(\mathcal{A})\quad\text{versus}\quad \lambda(\mathcal{L}).
\]
In particular, if $\lambda(\mathcal{A})\ge c\,\lambda(\mathcal{L})$ for some $c>1$, then \eqref{eq:tmix_gap_bound} implies
\[
t_{\mathrm{mix,\,levy}}^{(2)}(\delta)\le c^{-1}t_{\mathrm{mix,\,diff}}^{(2)}(\delta).
\]
\end{enumerate}

%--------------------------------------------------------



%--------------------------------------------------------
% \paragraph{Specialization to the L\'evy-score sampler.}
% Under Assumptions \ref{assumption} and \ref{ass:finite_activity_dirichlet},
% Proposition~\ref{thm:closed_form_energy} provides the explicit decomposition \eqref{eq:closed_form_energy}, which can be
% inserted into \eqref{eq:spectral_gap_def}--\eqref{eq:poincare_inequality} to estimate/lower bound $\lambda(\mathcal{A})$
% in concrete examples.

\begin{remark}
The $L^2(p_\infty)$ approach controls a $\chi^2$-type distance and is often convenient for spectral comparisons.
Other distances (e.g.\ total variation or Wasserstein) require different tools (e.g.\ minorization/coupling or contractive
couplings), and a spectral gap alone does not automatically give sharp bounds in those metrics without additional
assumptions.
\end{remark}


